
\section{Preliminary}

In this section, we first present a general formulation of a commonly used scaling law, and then demonstrate how to modify it to incorporate the vocabulary.

\subsection{Scaling law}

Scaling laws consider a computational budget, $C$, which is measured in FLOPs. The goal is to optimally allocate the compute budget to model parameters $N$ and the number of training tokens $D$ \citep{kaplan2020scaling,bahri2021explaining,hoffmann2022training,muennighoff2024scaling}. It can be formulated as:
\begin{equation}
    (N^{\rm opt}, D^{\rm opt}) = \arg\min_{N,D} \mathcal{L}(N,D) \quad \text{s.t. FLOPs}(N,D) = C,
\end{equation}
Following \citet{radford2018improving}, the loss function is typically the language modeling loss when evaluating language models, which can be written as:
\begin{equation}
 \mathcal{L}= -\frac{1}{T} \sum_{i=1}^{T} \log p(w_i | w_{1:i-1},V),
 \label{eq:lm_loss}
\end{equation}
where $p(w_i | w_{1:i-1},V)$ is the output probability of word $w_i$ given the context $w_{1:i-1}$ and the tokenizer with vocabulary size $V$.
Generally, the lower $\mathcal{L}$ indicates better performance of the language model.
However, due to its dependency on $V$, $\mathcal{L}$ cannot be used to compare language models with different vocabulary sizes. Thus, we propose an adaptation later in \autoref{sec:preliminaries}. Fitting scaling laws generally requires various models trained for different configurations~\citep{gadre2024language}. A common approach is to select several compute budgets and train models with varying $N$ and $D$ for each budget to find the best one, i.e. the one with the lowest loss (``IsoFLOPs'')~\citep{hoffmann2022training}. Using fitting techniques we can then estimate a function that maps from the compute budget to the optimal allocation to $N$ and $D$.



